\sezione{5}{Parametrico vs non parametrico}{23–25}

Sezione breve ma di snodo: spiega \emph{perché} il resto della lezione prende la strada che prende.

\subsection{Il problema di fondo}
\begin{center}
\rowcolors{2}{chiaro}{white}
\begin{tabularx}{\linewidth}{l l Y}
\rowcolor{navy}\thd{Ingrediente} & \thd{Difficoltà} & \thd{Perché}\\
Prior $P(w_i)$ & \textcolor{verde}{\textbf{facile}} & pochi numeri: conteggi, oppure classi equiprobabili\\
Likelihood $p(\x|w_i)$ & \textcolor{rosso}{\textbf{difficile}} & è un'intera \emph{funzione} su tutto lo spazio, da ricostruire a partire da pochi punti\\
\end{tabularx}
\end{center}
\begin{intuizione}
Stimare un prior è come stimare «la percentuale di maschi»: un solo numero. Stimare $p(\x|\text{maschi})$ è come disegnare la mappa completa di «dove stanno i maschi» nel piano peso-altezza, con la densità in \emph{ogni} punto, avendo visto solo 8 persone. Per questo le densità condizionali sono note «solo in teoria».
\end{intuizione}

\subsection{Le due soluzioni}
\begin{tcbraster}[raster columns=2,raster equal height,raster column skip=10pt]
\begin{tcolorbox}[base,colback=blu!6,colbacktitle=blu,coltitle=white,borderline west={3pt}{0pt}{blu},title={Parametrico},top=13pt,unbreakable]
«Scommetto che i dati hanno una certa \textbf{forma}; devo solo trovare i pochi numeri che la descrivono.»
\begin{itemize}
  \item si \emph{assume} la famiglia: gaussiana (multinormale);
  \item si stimano solo i \textbf{parametri} $\mmu$ e $\SSigma$;
  \item poi i dati di training si possono buttare.
\end{itemize}
\end{tcolorbox}
\begin{tcolorbox}[base,colback=ambra!7,colbacktitle=ambra,coltitle=white,borderline west={3pt}{0pt}{ambra},title={Non parametrico},top=13pt,unbreakable]
«Non assumo nulla sulla forma: lascio che siano i \textbf{dati} a disegnarla.»
\begin{itemize}
  \item nessuna forma imposta;
  \item la densità si costruisce dai punti (\emph{Parzen Window}; il cerchio della sezione 4 ne è una versione grezza);
  \item bisogna conservare tutto il training set.
\end{itemize}
\end{tcolorbox}
\end{tcbraster}

\begin{attenzione}[Il nome inganna]
«Non parametrico» \textbf{non} vuol dire «senza parametri». Vuol dire che il numero di parametri \textbf{cresce con i dati}: in un k-NN i dati stessi sono i parametri. Un modello parametrico ha invece un numero di parametri \emph{fisso}, qualunque sia la dimensione del training set.
\end{attenzione}

\subsection{Quando usare quale}
Il parametrico si usa quando valgono \textbf{entrambe} le condizioni: (1) c'è una «ragionevole certezza (o \emph{speranza})» che la forma sia adeguata, e (2) il training set è troppo piccolo per una buona stima non parametrica. Il corsivo su \emph{speranza} è voluto: spesso è un'ipotesi di comodo, da verificare (sezione 12).

\begin{center}
\rowcolors{2}{chiaro}{white}
\begin{tabularx}{\linewidth}{l Y Y}
\rowcolor{navy} & \thd{Parametrico} & \thd{Non parametrico}\\
Ipotesi sulla forma & forte (es. gaussiana) & nessuna\\
N. parametri & fisso & cresce con i dati\\
Dati necessari & pochi & molti\\
Rischio principale & \textbf{underfitting} se la forma è sbagliata & \textbf{overfitting} se i dati sono pochi\\
Memoria / velocità & leggero, veloce & conserva tutto, lento in predizione\\
Esempi & Bayes con multinormali, Naive Bayes, regressione logistica & Parzen Window, k-NN, alberi di decisione\\
\end{tabularx}
\end{center}

\subsection{Gradi di libertà e overfitting}
\begin{tecnica}[Gradi di libertà]
Numero di parametri indipendenti che il modello può regolare per adattarsi ai dati: misurano la sua \textbf{flessibilità} (capacità espressiva).
\begin{itemize}
  \item gaussiana 1D: \textbf{2} parametri ($\mu,\sigma$);
  \item multinormale in $d$ dimensioni: $d+\frac{d(d+1)}{2}$ per classe, cioè \textbf{5} con $d=2$;
  \item non parametrico: il modello «ricorda» ogni punto, quindi la flessibilità cresce con $n$.
\end{itemize}
\end{tecnica}
\begin{intuizione}[Perché pochi gradi di libertà = meno overfitting]
Un modello molto flessibile, con pochi dati, ha «troppa libertà»: si piega attorno a ogni punto, compresi quelli anomali (\textbf{overfitting}: va benissimo sul training e male sui dati nuovi). Un modello rigido come la gaussiana, con i suoi 5 numeri, \emph{non può} inseguire ogni punto: è costretto a catturare la tendenza generale. Il prezzo: se la forma è sbagliata non si adatta neanche con tanti dati (\textbf{underfitting}).
\end{intuizione}

\fig[.98\linewidth]{param_nonparam.pdf}
\didascalia{Sinistra: dati gaussiani e pochi campioni, la gaussiana stimata è già buona mentre le finestre producono picchi. Destra: dati bimodali e molti campioni, la gaussiana è sbagliata qualunque sia $n$, le finestre ricostruiscono entrambe le gobbe.}

\begin{mlbox}[Il compromesso bias-varianza]
Parametrico = \emph{bias} alto (se la forma è sbagliata), varianza bassa. Non parametrico = bias basso, \emph{varianza} alta con pochi dati. L'obiettivo di tutto il ML, come dice la slide 25: imparare i pattern reali nei dati \emph{senza memorizzare il rumore}. Anche il non parametrico ha un parametro da scegliere bene: la larghezza $h$ della finestra.
\end{mlbox}

\begin{appunti}{— Sezione 5}
\pag{23} Prior facili; likelihood difficile (intera funzione). \textbf{Parametrico}: si assume la forma, si stimano $\mmu,\SSigma$. \textbf{Non parametrico}: densità dai punti (Parzen), si conserva il training set. «Non parametrico» = n. parametri che cresce con i dati.
\pag{25} Parametrico se forma plausibile \textbf{e} pochi dati. Gradi di libertà = flessibilità: 2 (gaussiana 1D), $d+d(d+1)/2$ (multinormale). Pochi gradi $\Rightarrow$ meno overfitting; forma sbagliata $\Rightarrow$ underfitting. Compromesso bias-varianza.
\end{appunti}
